\section{Java effect engines: uplift, graph-compatible slopes and decision value}
\label{exec20:javafx-scope}

A servicing team may wish to rank customers by the incremental benefit of an
offer, while a treasury analyst may wish to rank possible causes of a funding
response. Both tasks can return a number called an effect, but the corresponding
software contracts differ. H2O-3 supplies a Java implementation of uplift
forests with separate treated and control predictions. Tetrad supplies
graph-compatible regression effects and a stability-ranking construction.
The proposed common library should preserve these useful distinctions rather
than turn every numeric output into a causal confidence interval. The selected
H2O source declares version 3.47.0 and a bleeding-edge profile at the pinned
20 September 2026 revision. Tetrad is inspected at its pinned 22 September 2026
revision. These identities specify the primary text being compared; they do
not certify a released binary or a completed integration.
\citep{exec20_javaeffects_h2o,exec20_javaeffects_tetrad}

\subsection{The target precedes the ranking algorithm}
\label{exec20:javafx-target}

Let $A\in\{0,1\}$ denote a servicing action, $Y(a)$ a binary desirable
response by a declared horizon, and $X$ the information available before the
action. The conditional effect of interest is
\begin{equation}
 \tau(x)=E[Y(1)-Y(0)\mid X=x].
 \label{exec20:javafx-cate}
\end{equation}
For the proposed public comparison, assume consistency, no interference,
conditional exchangeability of the two potential responses and $A$ given $X$,
and positive probability of both actions in the relevant covariate region.
Then $E[Y(a)\mid X=x]=E[Y\mid A=a,X=x]$, so the difference of the
two conditional response means identifies (\ref{exec20:javafx-cate}). This
follows by replacing the conditioning on $A$ under exchangeability and then
using consistency; a treatment-column name alone supplies neither step.
The argument is conditional on these imposed assumptions, not an assertion
that a bank's historical assignments satisfy them.
\citep{exec20_javaeffects_reference}

An exact example makes the cost of omitting assignment information visible.
Suppose half the customers are in each of two pre-action risk groups. Their
action probabilities are $1/5$ and $4/5$, and their response probabilities are
also $1/5$ and $4/5$. Within each group impose identical potential responses,
$Y(1)=Y(0)$, so the causal effect is zero. If the group is omitted,
\begin{equation}
 E[Y\mid A=1]=\frac{17}{25},\qquad
 E[Y\mid A=0]=\frac{8}{25},\qquad
 E[Y\mid A=1]-E[Y\mid A=0]=\frac9{25}.
 \label{exec20:javafx-confounding}
\end{equation}
For example, the first conditional mean is
$[\tfrac12(\tfrac15)^2+\tfrac12(\tfrac45)^2]/(\tfrac12)=17/25$.
With an assumed USD 100 value per incremental response and USD 20 action cost,
the unadjusted contrast suggests USD 16, while the stipulated causal margin is
USD $-20$. This construction does not show that an uplift forest fails when
the confounder is supplied and the model is adequate. It shows why assignment,
covariates, eligibility and target population must accompany a proposed engine.
\citep{exec20_javaeffects_reference}

\subsection{H2O's Java training and scoring contract}
\label{exec20:javafx-uplift}

The selected H2O builder accepts binomial response modeling and requires a
categorical treatment column with two levels, no missing treatment values and
internal extrema zero and one. These checks do not by themselves bind a
business label such as ``contacted'' to the intended treatment code. The
wrapper must save the category domain and the desirable response event.
The selected training path explicitly rejects weights, offsets, a fold column
and a positive built-in cross-validation fold count. A proposed weighted
observational analysis therefore cannot silently pass its weights through this
interface. External held-out evaluation is a separate possible design, subject
to its own sampling and training separation.
\citep{exec20_javaeffects_h2o}

The uplift tree uses a common split structure for both arms and retains two
leaf predictions. For an ordinary leaf with $n_1,n_0$ treated and control
observations and $s_1,s_0$ positive responses, the documented smoothing is
\begin{equation}
 \widehat p_1=\frac{s_1+1}{n_1+2},\qquad
 \widehat p_0=\frac{s_0+1}{n_0+2},\qquad
 \widehat u=\widehat p_1-\widehat p_0.
 \label{exec20:javafx-smoothing}
\end{equation}
The inspected ordinary split code constructs these smoothed rates; the forest
and the selected MOJO scorer average the two probability sums over trees and
return their difference. The three prediction entries are the uplift score,
the treated positive-response prediction and its control counterpart. A
probability difference is not the probability that this particular customer
benefits: that individual statement additionally depends on the joint law of
the two potential responses.
\citep{exec20_javaeffects_h2o}

For ten observations in each arm with eight versus two positive responses,
the raw contrast is $3/5$, but (\ref{exec20:javafx-smoothing}) gives
$3/4-1/4=1/2$. At an assumed USD 100 response value and USD 40 cost, these
two point calculations give USD 20 and USD 10. The difference is an explicit
consequence of smoothing, not empirical evidence of which estimate is better.
If the control arm has no observations, the same formula supplies $1/2$ for
its probability. That numerical value does not identify an unobserved control
response. The result contract should report the missing support alongside any
regularized prediction.
\citep{exec20_javaeffects_reference}

Missing-feature routing and a tree that never splits require additional
qualification. In the selected nonsplitting-root branch, the two leaf values
are taken from the two entries of \texttt{priorClassDist}. The route and
meaning of that inherited object must be traced before treating this branch
as an arm-specific estimate. The selected upstream depth-zero test checks that
a model exists; it does not resolve that semantic question. Other inspected
tests do assert prediction-column names and Java/MOJO scoring comparisons,
which are useful acceptance designs but have not been executed in this work.
\citep{exec20_javaeffects_h2o}

\subsection{A split criterion needs an executable mathematical identity}
\label{exec20:javafx-split}

H2O documents KL, squared Euclidean and chi-squared divergences for uplift
splits. The selected default resolves \texttt{AUTO} to KL, and the histogram
constructor selects the corresponding divergence object. The inspected KL
method contains the expression
\begin{quote}
\ttfamily\small
p * log2(p / q == 0 ? epsilon : q).
\end{quote}
Here $p$ and $q$ abbreviate its two probability arguments. The division and
comparison determine the conditional branch passed to the logarithm; the
expression is not the documented $p\log(p/q)$. For $p=q=1/2$, each selected
literal contribution is $-1/2$, so the two-class node sum is $-1$. The
documented KL sum is zero. The discrepancy persists regardless of the
logarithm base used for the documented zero.
\citep{exec20_javaeffects_h2o,exec20_javaeffects_reference}

This is a bounded source-and-arithmetic finding. It is not a report of an
executed Java failure, a statement about all H2O versions, or a measurement of
the resulting fitted forest. The next public prototype should evaluate the
selected method, normalization and actual split choices on equal probabilities,
zero-support boundaries and asymmetric arm counts. A literal port would carry
the discrepancy; replacing the formula without preserving the original would
create a new algorithm that needs separate review. An alternative divergence
is a candidate to compare, not an already-qualified repair.
\citep{exec20_javaeffects_h2o,exec20_javaeffects_reference}

\subsection{Average predictions and AUUC answer different questions}
\label{exec20:javafx-metrics}

The selected metric builder accumulates predicted uplift values and reports
their averages over all retained rows, treated rows and control rows as ATE,
ATT and ATC. It skips rows with a missing outcome. Thus both the prediction
function and the retained-row population matter; the labels do not by
themselves establish an independently identified or unbiased causal estimate.
For fixed scores $(0.1,0.2,0.3,0.4)$ with the first two rows treated, the
three averages are $1/4,3/20,7/20$. Excluding the fourth row because its
outcome is missing changes the all-row and control averages to $1/5$ and
$3/10$, while the treated average remains $3/20$. A production result must
explain such changes in its population rather than merely label them ATE.
\citep{exec20_javaeffects_h2o,exec20_javaeffects_reference}

For a prefix of customers ordered by a score, let $N_1,N_0$ be arm counts
and $S_1,S_0$ positive-response counts. The selected AUUC implementation uses
\begin{equation}
 L=\frac{S_1}{N_1}-\frac{S_0}{N_0},\qquad
 Q=S_1-S_0\frac{N_1}{N_0},\qquad
 G=(N_1+N_0)L,
 \label{exec20:javafx-curves}
\end{equation}
for lift, qini and gain. These expressions require support in the relevant
arms. Quantile thresholds determine bins; its area convention is a sum of
the cumulative curve value times the bin frequency, divided by $n+1$.
This convention must be fixed in a cross-engine comparison. The selected
implementation replaces a sole NaN curve value with zero and otherwise calls
an interpolation helper for NaNs. Such display processing does not create
the missing assignment information.
\citep{exec20_javaeffects_h2o}

Consider four ordered bins, each containing one treated and one control
customer, with observed response pairs $(1,0),(0,0),(1,1),(0,1)$. Their
cumulative qini values are $1,1,1,0$; the stated rectangle convention gives
$2(1+1+1+0)/9=2/3$, and the endpoint random-reference curve is zero. Scores
$(0.4,0.3,0.2,0.1)$ and their doubled values produce exactly the same order
and these same bins. Their average predictions are $1/4$ and $1/2$. Thus a
matched ranking curve does not establish probability-difference calibration.
This fixed-label arithmetic does not assert any potential-outcome law for
the eight customers. A causal policy comparison needs an independent design,
assignment support, a suitable value estimator and uncertainty for its sampling
unit. Reusing a validation sample repeatedly to choose a threshold also needs
to be included in that evaluation design.
\citep{exec20_javaeffects_reference}

\subsection{Tetrad's effects retain direction, graph ambiguity and failure}
\label{exec20:javafx-graphs}

Tetrad's selected IDA interface works with continuous data and legal partially
directed acyclic graphs. It enumerates admissible local parent choices and
computes regression coefficients using covariance submatrices, with indices
matched by variable names. This preserves a set of candidate linear slopes,
not a general nonlinear heterogeneous response surface. The unchanged IDA
source also underlies the earlier identification discussion; its reuse here
does not add an independent algorithmic coverage result. A selected singular
inverse catch returns zero, while other candidate exceptions can be logged
and omitted. These branches motivate a typed distinction among a graph-implied
zero, a numerical failure and an incomplete candidate collection.
\citep{exec20_javaeffects_tetrad}

For example, if two regressors are identical, a covariance system can determine
only $\beta_X+\beta_Z=1$. The solutions $(0,1),(1,0),(-1,2)$ give the
same fitted combination. Returning zero for the first coefficient does not
show that its causal effect is identified as zero. This exact algebra is a
proposed numerical qualification case; it has not been run through Tetrad.
\citep{exec20_javaeffects_reference}

CStaR extends the comparison from a graph effect set to a magnitude ranking
across subsamples. Its selected code takes the minimum of the absolute
candidate effects. For candidates $\{-2,-1/2\}$, this is $1/2$, whereas
taking the absolute value of the signed minimum gives $2$. The first is the
inspected implementation's intended magnitude summary, but neither operation
retains the sign. At a hypothetical USD 100 value per response unit and USD 20
action cost, the two signed candidates imply margins $-220$ and $-70$.
Misreading the minimum magnitude as a positive benefit gives USD 30. An
integration should retain the signed candidates and their units beside the
ranking summary.
\citep{exec20_javaeffects_tetrad,exec20_javaeffects_reference}

The selected default uses thirty half-samples without replacement; its
bootstrap option also draws half the rows, with replacement. Neither route
alone specifies branch-cluster or shared-calendar resampling. The ranking
uses an inclusive cutoff at the requested bracket times the number of effect
variables. With one effect, bracket one, and three equal positive magnitudes,
all three meet the cutoff. If this happens in every subsample, each has
selection frequency one. That frequency is a property of the construction;
it is not automatically a probability of causation or an interval's coverage.
The displayed PCER calculation uses the displayed rank as its $q$ argument.
Its presence does not verify the theoretical assumptions behind an error bound
for the proposed data. The source also states that required and forbidden
edge knowledge is not respected by this class. A business constraint must
therefore be preserved through a verified route or explicitly refused.
\citep{exec20_javaeffects_tetrad,exec20_javaeffects_reference}

\subsection{A proposed adoption route and a public qualification experiment}
\label{exec20:javafx-adoption}

H2O is a candidate for an original Java training worker with a separate MOJO
scoring service; Tetrad is a candidate for graph-effect and stability-ranking
reference services. The two should expose different result types through the
library, CLI and MCP surfaces. A request should bind row and variable identities,
action and positive-event mappings, horizon, eligible population, graph type,
original units, sampling law, source revision and executable artifact hash.
A response should separate prediction, causal estimand, graph candidates,
ranking metric, numerical status and qualified uncertainty. None of these
proposed Java services has been implemented or admitted by this source review.
\citep{exec20_javaeffects_h2o,exec20_javaeffects_tetrad}

Preserving the original engines is preferable as the first comparison because
it makes a source-to-output discrepancy observable before maintaining a port.
H2O's training and generated-model modules have different dependencies; the
generated-model build's compatibility setting is not a statement about every
training deployment. Its selected MOJO subclasses both declare format 1.40,
but use different threshold-key strings in their writer and reader. Superclass
and default handling must be traced and a full archive roundtrip tested before
concluding whether this affects scoring or metadata. The root Apache 2.0
declaration and Tetrad's GPLv3-or-later source headers are relevant inputs to
reuse review; neither a process boundary nor a source notice resolves the
complete distribution and dependency obligations.
\citep{exec20_javaeffects_h2o,exec20_javaeffects_tetrad}

The first public experiment should therefore use frozen synthetic mechanisms
with known action effects, followed by appropriately documented public data.
The synthetic suite should include the confounded zero-effect law, missing arm
support, constant/root-only trees, missing features, category and column
permutations, the equal-probability split witness, signed graph effects,
singular covariance and ranking ties. Separate checks should match engine and
MOJO predictions and preserve keyed data through each transport. After that
semantic comparison, a fresh held-out study can evaluate calibration, ranking,
policy value and repeated-fit costs under declared independent or dependent
sampling laws. No bank experiment is required to establish whether these
public prototype steps merit the next investment.
\citep{exec20_javaeffects_reference}

This review supplies two concrete Java profiles rather than an absence finding
from an empty search. Saved complementary repository leads and search remainder
still need original-source inspection. Foundational papers named in upstream
comments need genuine local PDFs and technical review before their theorems
are attributed here. Those tasks can proceed alongside the pinned public
prototype design; their openness does not justify converting a score, a
frequency or a successful API call into an unsupported causal guarantee.
\citep{exec20_javaeffects_h2o,exec20_javaeffects_tetrad}
